\section{Notions et définitions}
\label{annexes:notions}

\subsection{Forme conservative}
\label{annexes:forme_conservative}

Une équation sous forme conservative exprime un bilan local :
la variation d'une quantité par unité de volume résulte des flux
qui la transportent et d'éventuelles sources.
Pour une grandeur scalaire $q$, elle s'écrit :
\begin{equation*}
    \frac{\partial q}{\partial t} + \nabla \cdot \mathbf{F} = s,
\end{equation*}
où $\mathbf{F}$ est le flux et $s$ le terme source.
Sans source ni flux à travers les frontières, la quantité totale
dans un domaine fixe est conservée.
Par exemple, le bilan de masse utilise $q = \rho$,
$\mathbf{F} = \rho\mathbf{u}$ et $s = 0$.

\subsection{Tenseur des contraintes et modèle newtonien isotrope}
\label{annexes:type_fluids}

Le tenseur des contraintes de Cauchy décrit les efforts internes
et se décompose ici en une contribution de pression et une contribution
visqueuse : $\boldsymbol{\sigma} = -p\mathbf{I} + \boldsymbol{\tau}$.
Le modèle newtonien isotrope suppose une relation linéaire entre
les contraintes visqueuses et le tenseur des taux de déformation,
sans direction matérielle privilégiée.
La dérivation de cette relation et les hypothèses associées seront
détaillées ultérieurement.
